\section{On Circuits and Restrictions}
\label{sec:circ-rest}

Let $G=(U, V, E)$ be a bipartite graph with $U = \set{0,1}^n$ and
$V = [m]$. As in the XOR-CSP setup (\cref{sec:known-SoS-results}) we
think of vertices in $U$ as constraints and vertices in $V$ as
variables. More specifically, we think of each vertex $\alpha \in U$
as an \emph{xor} constraint over the variables in the neighborhood
$\oplus_{i \in N(\alpha)} v_i = b_\alpha$, for a constraint vector
$b \in \set{0,1}^U$. Given an assignment $\beta \in \set{0,1}^m$ to
the variables $V$, we let $f_{G, \beta}: U \rightarrow \set{0,1}$ be
the function defined by
$f_{G, \beta}(\alpha) = \oplus_{i \in N(\alpha)} v_i$.  In other
words, viewing $f_{G,\beta}$ as a vector in $\set{0,1}^U$, it is the
unique constraint vector such that the XOR-CSP instance, defined over
$G$, is satisfied by the assignment $\beta$. Let us denote the set of
all such constraint vectors that give rise to a satisfiable XOR-CSP
instance by
\[
  \calF_G = \set{f_{G, \beta} \,|\, \beta \in \set{0,1}^m} \eqperiod
\]
In order for SoS to refute an XOR-CSP instance defined over $G$,
it must prove that the given constraint vector is not in the set
$\calF_G$.

On the other hand in order for SoS to refute the formula $\circuit_s(f)$
it needs to show that there is no circuit of size at most $s$
computing $f$. That is, SoS needs to show that $f$ is not in the set
\[
  \calC_{\emptyset} =
  \set{
    T:
    \set{0,1}^n \rightarrow \set{0,1}
    \text{~such that~}
    \circuit_s(T)
    \text{~is satisfiable}
  } \eqperiod
\]
More generally, if we restrict $\circuit_s(f)$ by a restriction
$\rho$, then the proof system must prove that $f$ is not a member
of the family of truthtables
\[
  \calC_\rho =
  \set{
    T:
    \set{0,1}^n \rightarrow \set{0,1}
    \text{~such that~}
    \restrict{\circuit_s(T)}{\rho}
    \text{~is satisfiable}
  } \eqperiod
\]

In the following we show that there is a well-behaved restriction
$\rho$ such that $\calC_\rho = \calF_G$ for some explicit graphs
$G$. In other words, once we consider the restricted formula
$\restrict{\circuit_s(f)}{\rho}$, SoS needs to rule
out that $f$ is a valid right hand side of an XOR-CSP instance. But we
know that if $G$ is a moderate expander, then low degree SoS cannot
determine whether the XOR-CSP instance is satisfiable and hence we
obtain our lower bound.

Let us first formalize the properties we require from $\rho$. We start
off by restricting our attention to a certain natural class of variable
substitutions. Namely, we do not want that the structure of the circuit
depends on evaluation variables.

\begin{definition}[natural variable substitutions]
  A variable substitution $\rho$ to the variables of $\circuit_s(f)$ is
  \emph{natural} if there is \emph{no} structure variable $y$ such
  that $\rho(y)$ is an evaluation variable.
\end{definition}

In order to motivate the next definition, let us informally
describe the natural restriction $\rho$ and explain the properties of
$\rho$ we require.

For now we can think of $\rho$ as a restriction to the structure
variables (though for the size lower bounds we also need to restrict
some of the evaluation variables). Some set of $m$ structure variables remains
undetermined. Let us denote these variables by $y_1, \ldots, y_m$. We
intend to choose $\rho$ such that on a given input
$\alpha \in \set{0,1}^n$ to the circuit, it is forced to compute
$\oplus_{i \in N(\alpha)} y_i$.  In other words, given such a
restriction $\rho$, we are \emph{essentially} left with an XOR-CSP
problem over $G$, with right hand side $f$.  There is however a
difference in that the encoding is non-standard: the evaluation
variables act like extension variables that correspond to the
functions computed at each gate of the circuit. In order to argue that
the known degree lower bound for the XOR-CSP problem implies a degree
lower bound for the problem at hand, we need to get rid of these
extension variables. This can be done if the functions computed at the
gates are of low degree in the $y$ variables.

Recall from \cref{sec:restrictions} that a system of polynomial
equations $\calP$ has $n$ unset variables if there are $n$ tuples of
variables $(x, \bar x)$ such that at least one variable of each tuple
occurs in $\calP$ and all variables in these tuples are unset, i.e.,
they are not fixed to a constant.

\begin{definition}[$k$-determined]
  \label{def:determined}%
  Let $\rho$ be a variable substitution to the variables of
  $\circuit_s(f)$ and suppose that $\rho$ leaves $m$ structural
  variables $Y = \set{y_1, \ldots, y_m}$ unset.  Then $\rho$ is
  \emph{$k$-determined} if for every $v \in [s]$ and
  $\alpha \in \set{0,1}^n$ there are multilinear polynomials
  $$\gout_{v, \alpha}, \ginone_{v, \alpha}, \gintwo_{v, \alpha}:
  \set{0,1}^m \rightarrow \set{0,1}$$ depending on at most $k$
  variables such that the following holds. For all $T \in \calC_\rho$
  and all total assignments $\sigma$ that satisfy
  $\restrict{\circuit_s(T)}{\rho}$ it holds that
  \begin{align}
    \label{eq:k-determined-substitution}
    \restrict{\outwire_\alpha(v)}{\rho \cup \sigma} &=
    \gout_{v, \alpha}(\beta) \eqcomma &
    \restrict{\inwire_\alpha(v, 1)}{\rho \cup \sigma} &= \ginone_{v,
    \alpha}(\beta) \eqcomma \text{~and} &
    \restrict{\inwire_\alpha(v, 2)}{\rho \cup \sigma} &= \gintwo_{v,
    \alpha}(\beta) \eqcomma
  \end{align}
  where $\beta \subseteq \sigma$ is the assignment to $Y$.
\end{definition}

However, \cref{def:determined} is not quite sufficient. For example,
there is no guarantee that $\calC_\rho$ is non-empty, i.e., that the
restriction $\rho$ describes a valid (partial) circuit. More generally, we
need the additional guarantee that there are still many viable
circuits that the restricted formula can describe: if there is just a
single setting of the $Y$ variables such that all structural axioms
are satisfied, then the formula may be refuted in constant
degree. Hence we need to ensure that there are many viable assignments
to the $Y$ variables that satisfy all structure axioms. This leads us
to the following definition.

\begin{definition}[$m$-independent]
  A variable substitution $\rho$ to the variables of the formula
  $\circuit_s(f)$ is \emph{$m$-independent} if $\rho$ leaves exactly
  $m$ structural variables $Y = \set{y_1, \ldots, y_m}$ unset, and for
  every assignment $\beta \in \set{0,1}^Y$ it holds that
  $|\calC_{\rho \cup \beta}| = 1$.
\end{definition}

With these definitions at hand we can state the lemma that drives all
our lower bounds.

\begin{lemma}
  \label{lem:circuit-to-poly}
  Let $\rho$ be a natural $m$-independent $k$-determined variable substitution
  of $\circuit_s(f)$, and let $Y$ and
  $\gout_{u, \alpha}$ be as in \cref{def:determined}. If there is an
  SoS refutation of $\restrict{\circuit_s(f)}{\rho}$ of degree $d$,
  then there is a degree $d \cdot k$ SoS refutation of the system
  of polynomial equations
  \begin{align}
    \label{eqn:lb transfer lemma}
    \set{
      \gout_{s,\alpha}(Y) = f(\alpha) \mid \alpha \in \set{0,1}^n
    }
    \cup
    \set{
      y_i^2 = y_i \mid i \in [m]
    } \eqperiod
  \end{align}
\end{lemma}


For proving this lemma, we consider the natural extension of $\rho$
which substitutes all evaluation variables by appropriate
degree-$k$ polynomials as indicated by \cref{def:determined}.

\begin{definition}
  For a $k$-determined restriction $\rho$ (with associated polynomials
  $\gout_{v,\alpha}, \ginone_{v,\alpha}, \gintwo_{v,\alpha}$) of
  $\circuit_s(f)$, we denote by $\hat\rho$ the polynomial substitution
  that extends $\rho$ by first substituting any bar variable $\bar x$
  by $1 - x$ and then substituting all evaluation variables as
  follows:
  \begin{align*}
    \hat\rho(\outwire_\alpha(v)) &=
    \gout_{v, \alpha}(Y) \eqcomma &
    \hat\rho(\inwire_\alpha(v, 1)) &= \ginone_{v,
    \alpha}(Y) \eqcomma \text{~and} &
    \hat\rho(\inwire_\alpha(v, 2)) &= \gintwo_{v,
    \alpha}(Y) \eqperiod
  \end{align*}
\end{definition}

Note that the formula $\restrict{\circuit_s(f)}{\hat\rho}$ is
defined only over $Y$. Let us stress that there are no ``bar''
variables left in the formula. The main observation used to prove
\cref{lem:circuit-to-poly} is the following claim, which establishes
that the formula \eqref{eqn:lb transfer lemma} is in fact essentially
the same as $\restrict{\circuit_s(f)}{\hat\rho}$.

\begin{claim}
  \label{claim:tau-rho}
  Let $\rho$ be a natural $m$-independent $k$-determined variable substitution
  of $\circuit_s(f)$.  Then
  $\restrict{\circuit_s(f)}{\hat\rho}$ can be written as
  $$
  \restrict{\circuit_s(f)}{\hat\rho} = \calP \cup \calQ \eqcomma
  $$
  where $\calP$ is the formula \eqref{eqn:lb transfer lemma} and
  $\calQ$ only consists of axioms that are satisfied for all
  assignments $\beta \in \{0,1\}^Y$.
\end{claim}

\begin{proof}
  Note that the set of output axioms \eqref{eq:ax-correct} of
  $\circuit_s(f)$ under $\hat\rho$ equals the first part of
  \eqref{eqn:lb transfer lemma}, and that the Boolean axioms on the $Y$ variables in
  $\restrict{\circuit_s(f)}{\hat\rho}$ are exactly the second part
  of \eqref{eqn:lb transfer lemma}.

  The remaining axioms of $\restrict{\circuit_s(f)}{\hat\rho}$, which
  are not present in \eqref{eqn:lb transfer lemma}, are the Boolean
  axioms on the variables outside $Y$, the negation axioms, as well as
  \cref{eq:ax-from,eq:ax-fn,eq:ax-in-gt,eq:ax-in-gt2,eq:prod-zero-in,eq:prod-zero-gate,eq:inwire-const,eq:inwire-var,eq:inwire-gate,eq:ax-neg,eq:ax-or,eq:ax-and}.

  The Boolean axioms may turn into polynomials of degree at most
  $2k$. Because the polynomials we substitute the variables with are
  Boolean valued, we see that these substituted axioms are satisfied
  for all assignments $\beta \in \set{0,1}^Y$ and we can thus put them
  into the set $\calQ$.

  The negation axioms all become ``$0=0$'' under $\hat{\rho}$ since
  $\hat{\rho}(\bar{x}) = 1 - \hat{\rho}(x)$.
  
  Finally we need to argue that the
  \cref{eq:ax-from,eq:ax-fn,eq:ax-in-gt,eq:ax-in-gt2,eq:prod-zero-in,eq:prod-zero-gate,eq:inwire-const,eq:inwire-var,eq:inwire-gate,eq:ax-neg,eq:ax-or,eq:ax-and}
  are also of the form $p(Y)=0$ for a polynomial $p$ which is
  identically $0$ on all of $\{0,1\}^m$.  This in turn follows
  immediately from the assumption that $\rho$ is $m$-independent: for
  every $\beta \in \{0,1\}^Y$, there exists some $T$ such that the
  complete assignment $\hat\rho(\beta) \cup \beta$ satisfies
  $\circuit_s(T)$.  But since none of the remaining
  \cref{eq:ax-from,eq:ax-fn,eq:ax-in-gt,eq:ax-in-gt2,eq:prod-zero-in,eq:prod-zero-gate,eq:inwire-const,eq:inwire-var,eq:inwire-gate,eq:ax-neg,eq:ax-or,eq:ax-and}
  depends on $T$, they must then all be satisfied for every
  $\beta \in \{0,1\}^Y$.
\end{proof}

Using this claim we can easily prove \cref{lem:circuit-to-poly}.

\begin{proof}[Proof of \cref{lem:circuit-to-poly}]
  Suppose $\restrict{\circuit_s(f)}{\rho}$ has a refutation in degree
  $d$.  By \cref{lem:sos degree polynomial substitution}, there then
  exists a degree $d \cdot k$ refutation of
  $\restrict{\circuit_s(f)}{\hat\rho}$.

  By \cref{claim:tau-rho}, this new refutation is \emph{almost} a
  refutation of \eqref{eqn:lb transfer lemma}, except that
  $\restrict{\circuit_s(f)}{\hat\rho}$ has an additional set $\calQ$
  of axioms that the refutation may use.  However, each of these
  additional axioms is of the form $p(Y) = 0$ for a polynomial which
  is identically $0$ on the entire Boolean cube.  By \cref{claim:0 mod
    bool}, such an axiom can be rewritten as a linear combination of
  the Boolean axioms.  Since the Boolean axioms are present in
  \eqref{eqn:lb transfer lemma}, this yields a refutation of that
  formula in degree $d \cdot k$.
\end{proof}


